% ============================================================
% 4_Experiments_and_Results.tex
% [待实验验证] = simulated estimate, replace after AutoDL runs
% ============================================================

\section{Experiments and Results}
\label{sec:experiments_results}

We organize experiments into seven series (E1--E7) totaling 63 controlled
comparisons. Table~\ref{tab:datasets} summarizes the three datasets used.
\textbf{All numerical values marked [待实验验证] are simulated estimates}
based on prior single-corpus results and will be replaced with measured values
from AutoDL execution.

\begin{table}[t]
\centering
\caption{Dataset statistics after preprocessing.}
\label{tab:datasets}
\begin{tabular}{lcccccc}
\toprule
\textbf{Dataset} & \textbf{Classes} & \textbf{Samples} & \textbf{Speakers} & \textbf{Age} & \textbf{Style} & \textbf{Language} \\
\midrule
C-BESD (MY) & 6 & 4,179 & 70 children & child (unk.) & acted & Malay \\
FAU Aibo    & 4 & 18,216 & 51 children & 10--13 & spontaneous & German \\
IEMOCAP     & 4 & 8,525 & 10 adults & adult & acted & English \\
\bottomrule
\end{tabular}
\end{table}

\subsection{E1: In-Domain Pooling Baseline}
\label{ssec:e1}

\textbf{Goal:} Establish the optimal pooling strategy for each dataset under
matched parameter budgets, serving as the unified baseline for all subsequent
experiments.

\begin{table}[t]
\centering
\caption{E1: In-domain pooling comparison (3 seeds, mean$\pm$std). [待实验验证]}
\label{tab:e1}
\begin{tabular}{llccc}
\toprule
\textbf{ID} & \textbf{Dataset} & \textbf{Pooling} & \textbf{Test WA (\%)} & \textbf{Test UAR (\%)} \\
\midrule
E1-01 & C-BESD (6cl) & Mean       & 78.0 $\pm$ 0.5 & 78.0 $\pm$ 0.5 \\
E1-02 & C-BESD (6cl) & Self-Attn  & \textbf{93.0 $\pm$ 1.5} & \textbf{93.0 $\pm$ 1.5} \\
E1-03 & C-BESD (6cl) & Prosody    & 91.0 $\pm$ 2.5 & 91.0 $\pm$ 2.5 \\
\midrule
E1-04 & FAU (4cl)    & Mean       & 62.0 $\pm$ 1.0 & 50.0 $\pm$ 2.0 \\
E1-05 & FAU (4cl)    & Self-Attn  & \textbf{66.5 $\pm$ 0.8} & \textbf{55.5 $\pm$ 2.0} \\
E1-06 & FAU (4cl)    & Prosody    & 65.0 $\pm$ 1.2 & 54.0 $\pm$ 2.2 \\
\midrule
E1-07 & IEMOCAP (4cl)& Mean       & 54.5 $\pm$ 0.5 & 42.0 $\pm$ 1.0 \\
E1-08 & IEMOCAP (4cl)& Self-Attn  & \textbf{76.0 $\pm$ 8.0} & \textbf{68.0 $\pm$ 5.0} \\
E1-09 & IEMOCAP (4cl)& Prosody    & 66.0 $\pm$ 7.5 & 58.0 $\pm$ 5.0 \\
\bottomrule
\end{tabular}
\end{table}

\noindent\textbf{Key observation [待实验验证]:} Self-Attention pooling achieves the
highest in-domain WA on all three datasets, but its advantage over Prosody
Guided pooling is marginal on FAU (+1.5pp) compared to C-BESD (+2.0pp) and
IEMOCAP (+10.0pp). The narrow gap on spontaneous child speech suggests prosody
priors partially compensate for environmental noise. We select Self-Attention
as the default pooling for all downstream experiments (E2--E7) unless otherwise
noted.

\subsection{E2: WavLM Frozen vs.\ Unfrozen}
\label{ssec:wavlm_unfreeze}

\textbf{Goal:} Verify that the frozen-backbone paradigm is necessary for
children's SER given limited data.

\begin{table}[t]
\centering
\caption{E2: Frozen vs.\ unfrozen WavLM (Self-Attn pooling, 3 seeds). [待实验验证]}
\label{tab:e2}
\begin{tabular}{llccc}
\toprule
\textbf{ID} & \textbf{Dataset} & \textbf{WavLM} & \textbf{Test WA (\%)} & \textbf{Trainable Params} \\
\midrule
E2-01 & C-BESD & Frozen   & \textbf{93.0 $\pm$ 1.5} & 704K \\
E2-01b& C-BESD & Unfrozen & 78.0 $\pm$ 3.0 & 95M \\
\midrule
E2-02 & FAU    & Frozen   & \textbf{66.5 $\pm$ 0.8} & 704K \\
E2-02b& FAU    & Unfrozen & 58.0 $\pm$ 5.0 & 95M \\
\midrule
E2-03 & IEMOCAP & Frozen  & 76.0 $\pm$ 8.0 & 704K \\
E2-03b& IEMOCAP & Unfrozen& \textbf{78.0 $\pm$ 5.0} & 95M \\
\bottomrule
\end{tabular}
\end{table}

\noindent\textbf{Key observation [待实验验证]:} Unfreezing WavLM severely degrades
performance on children's datasets (C-BESD: $-15$pp; FAU: $-8$pp) due to
overfitting on limited speaker diversity, while providing a marginal gain on
IEMOCAP (+2pp) where adult speech patterns benefit from fine-tuning. This
confirms the frozen-backbone paradigm as essential for child SER.

\subsection{E3: Zero-Shot Cross-Corpus Transfer Matrix}
\label{ssec:e3}

\textbf{Goal:} Measure the relationship between corpus-level distribution shift
and zero-shot transfer accuracy.

\begin{table}[t]
\centering
\caption{E3: Zero-shot transfer matrix (WA\%, Self-Attn pooling, seed 42). [待实验验证]}
\label{tab:e3}
\begin{tabular}{lccc|c}
\toprule
\textbf{Train $\downarrow$ / Test $\rightarrow$} & \textbf{C-BESD (4cl)} & \textbf{FAU} & \textbf{IEMOCAP} & \textbf{FD (train$\rightarrow$test)} \\
\midrule
C-BESD (4cl) & \textbf{93.0} & 20.0 & 33.0 & -- / 8.5 / 7.2 \\
FAU          & 25.5 & \textbf{66.5} & 24.5 & 8.5 / -- / ? \\
IEMOCAP      & 33.5 & 21.5 & \textbf{76.0} & 7.2 / ? / -- \\
\bottomrule
\end{tabular}
\end{table}

\noindent\textbf{Key observation [待实验验证]:} Zero-shot transfer fails uniformly
(all cross-corpus results $<$ 34\%), with the lowest performance on the
spontaneous target (FAU, 19.6--21.5\%). FD values computed between corpus pairs
show that larger distribution shifts predict lower transfer accuracy
(Spearman's $\rho \approx -0.9$ [待实验验证]).

\subsection{E4: Data Augmentation Sensitivity}
\label{ssec:e4}

\textbf{Goal:} Quantify how different augmentation strategies shift the data
distribution and impact accuracy, testing whether adult-derived augmentation
experience transfers to children.

\begin{table}[t]
\centering
\caption{E4: Augmentation sensitivity (Self-Attn pooling, 3 seeds). FD measured vs.\ clean training set. [待实验验证]}
\label{tab:e4}
\begin{tabular}{llccc}
\toprule
\textbf{ID} & \textbf{Dataset + Condition} & \textbf{FD} & \textbf{Test WA (\%)} & \textbf{$\Delta$WA} \\
\midrule
E4-01 & C-BESD Clean        & 0.00  & 93.0 $\pm$ 1.5 & -- \\
E4-02 & C-BESD + Adult data  & 9.87  & 68.0 $\pm$ 2.0 & $-$25.0 \\
E4-03 & C-BESD + AWGN (child)& 8.71  & 66.0 $\pm$ 2.5 & $-$27.0 \\
E4-04 & C-BESD + Extreme     & 12.0  & 46.5 $\pm$ 3.0 & $-$46.5 \\
\midrule
E4-05 & FAU Clean            & 0.00  & 66.5 $\pm$ 0.8 & -- \\
E4-06 & FAU + Adult data     & 5.00  & 52.0 $\pm$ 3.0 & $-$14.5 \\
E4-07 & FAU + AWGN           & 4.50  & 56.0 $\pm$ 2.0 & $-$10.5 \\
E4-08 & FAU + Extreme        & 8.00  & 42.0 $\pm$ 4.0 & $-$24.5 \\
\midrule
E4-09 & IEMOCAP Clean        & 0.00  & 76.0 $\pm$ 8.0 & -- \\
E4-10 & IEMOCAP + Child data & 6.87  & 68.0 $\pm$ 5.0 & $-$8.0 \\
E4-11 & IEMOCAP + AWGN       & 3.00  & 72.0 $\pm$ 6.0 & $-$4.0 \\
E4-12 & IEMOCAP + Extreme    & 9.00  & 54.0 $\pm$ 7.0 & $-$22.0 \\
\bottomrule
\end{tabular}
\end{table}

\noindent\textbf{Key observation [待实验验证]:} Adding adult speech to children's
training data is the most harmful augmentation for C-BESD ($-$25pp), confirming
that adult and child emotional speech occupy distinct regions of WavLM feature
space. The FD--WA relationship holds strictly: each increase in FD predicts a
decrease in WA, with stronger degradation when the augmentation source is
distributionally farther from the target.

\subsection{E5: Layer Fusion Ablation}
\label{ssec:layer_ablation}

\textbf{Goal:} Quantify the contribution of learnable 12-layer weighted fusion
over single-layer baselines.

\begin{table}[t]
\centering
\caption{E5: Layer fusion ablation (Self-Attn pooling, 3 seeds). [待实验验证]}
\label{tab:e5}
\begin{tabular}{llcc}
\toprule
\textbf{ID} & \textbf{Dataset} & \textbf{Fusion Method} & \textbf{Test WA (\%)} \\
\midrule
E5-01 & C-BESD & Last Layer (L12)       & 80.0 $\pm$ 2.0 \\
E5-02 & C-BESD & Best Single (grid)     & 84.0 $\pm$ 1.5 \\
E5-03 & C-BESD & \textbf{Weighted Sum}  & \textbf{93.0 $\pm$ 1.5} \\
\midrule
E5-04 & FAU    & Last Layer (L12)       & 52.0 $\pm$ 2.0 \\
E5-05 & FAU    & Best Single (grid)     & 58.0 $\pm$ 1.5 \\
E5-06 & FAU    & \textbf{Weighted Sum}  & \textbf{66.5 $\pm$ 0.8} \\
\midrule
E5-07 & IEMOCAP& Last Layer (L12)       & 62.0 $\pm$ 5.0 \\
E5-08 & IEMOCAP& Best Single (grid)     & 68.0 $\pm$ 4.0 \\
E5-09 & IEMOCAP& \textbf{Weighted Sum}  & \textbf{76.0 $\pm$ 8.0} \\
\bottomrule
\end{tabular}
\end{table}

\noindent\textbf{Key observation [待实验验证]:} Weighted sum fusion outperforms
both baselines by large margins across all datasets (+9--15pp over last layer,
+8--9pp over best single layer). The learned weight distributions consistently
favor mid-to-upper layers (L7--L9 receiving highest weights), with entropy
$\approx$2.48 indicating near-uniform distribution that still provides
meaningful gain through differential layer emphasis.

\subsection{E6: Module Ablation}
\label{ssec:e6}

\textbf{Goal:} Isolate the independent contribution of each architectural
component.

\begin{table}[t]
\centering
\caption{E6: Module ablation on C-BESD (6cl) and FAU (4cl), 3 seeds. [待实验验证]}
\label{tab:e6}
\begin{tabular}{lcccccc}
\toprule
\textbf{ID} & \textbf{Dataset} & \textbf{Adapter} & \textbf{Pooling} & \textbf{LayerFusion} & \textbf{WA (\%)} \\
\midrule
E6-01 & C-BESD & --    & Mean    & LastL   & 68.0 $\pm$ 1.0 \\
E6-02 & C-BESD & Yes   & Mean    & LastL   & 70.0 $\pm$ 1.0 \\
E6-03 & C-BESD & --    & SelfAttn& LastL   & 80.0 $\pm$ 2.0 \\
E6-04 & C-BESD & --    & SelfAttn& Weighted& 90.0 $\pm$ 2.0 \\
E6-05 & C-BESD & Yes   & SelfAttn& Weighted& \textbf{93.0 $\pm$ 1.5} \\
\midrule
E6-06 & FAU    & Yes   & SelfAttn& Weighted& \textbf{66.5 $\pm$ 0.8} \\
\bottomrule
\end{tabular}
\end{table}

\noindent\textbf{Key observation [待实验验证]:} Layer fusion contributes the
largest single-module gain (+10pp over last-layer baseline on C-BESD),
followed by Self-Attention pooling (+12pp over mean pooling). The Adapter
module provides only marginal benefit (+2pp), consistent with prior findings
for frozen-backbone SER.

\subsection{E7: Model Transfer (Fine-Tuning)}
\label{ssec:e7}

\textbf{Goal:} Test whether pre-training on one corpus and fine-tuning on
another improves over training from scratch on the target corpus.

\begin{table}[t]
\centering
\caption{E7: Model transfer (pre-train $\rightarrow$ fine-tune, 3 seeds). [待实验验证]}
\label{tab:e7}
\begin{tabular}{lllcc}
\toprule
\textbf{ID} & \textbf{Pre-train} & \textbf{Fine-tune} & \textbf{Target WA (\%)} & \textbf{Scratch WA} \\
\midrule
E7-01 & C-BESD & FAU    & 62.0 $\pm$ 1.5 & 66.5 \\
E7-02 & C-BESD & IEMOCAP& 45.0 $\pm$ 3.0 & 76.0 \\
E7-03 & FAU    & C-BESD & 80.0 $\pm$ 2.0 & 93.0 \\
E7-04 & FAU    & IEMOCAP& 38.0 $\pm$ 3.0 & 76.0 \\
E7-05 & IEMOCAP& C-BESD & 82.0 $\pm$ 2.0 & 93.0 \\
E7-06 & IEMOCAP& FAU    & 36.0 $\pm$ 3.0 & 66.5 \\
\bottomrule
\end{tabular}
\end{table}

\noindent\textbf{Key observation [待实验验证]:} Pre-training on a different corpus
consistently under-performs training from scratch on the target ($-$4 to
$-$38pp). The degradation is most severe when transferring from adult
(IEMOCAP) to children's data (FAU: $-$30pp; C-BESD: $-$11pp), reinforcing the
non-transferability of adult-derived representations.

% ============================================================
% Summary table placeholder
% ============================================================
\subsection{Summary of Experimental Results}
\label{ssec:summary}

Table~\ref{tab:summary} provides a consolidated view of all 63 experiments.
\textbf{All values [待实验验证] will be updated after AutoDL execution.}

\begin{table}[p]
\centering
\caption{Complete experimental results summary. Bold = best per group. $\dagger$ = single seed. [待实验验证]}
\label{tab:summary}
\small
\begin{tabular}{llllcc}
\toprule
\textbf{Series} & \textbf{ID} & \textbf{Setting} & \textbf{Pooling/Condition} & \textbf{WA (\%)} & \textbf{UAR (\%)} \\
\midrule
\multirow{9}{*}{E1}
& E1-01 & C-BESD 6cl & Mean & 78.0$\pm$0.5 & 78.0$\pm$0.5 \\
& E1-02 & C-BESD 6cl & Self-Attn & \textbf{93.0$\pm$1.5} & \textbf{93.0$\pm$1.5} \\
& E1-03 & C-BESD 6cl & Prosody & 91.0$\pm$2.5 & 91.0$\pm$2.5 \\
& E1-04 & FAU 4cl & Mean & 62.0$\pm$1.0 & 50.0$\pm$2.0 \\
& E1-05 & FAU 4cl & Self-Attn & \textbf{66.5$\pm$0.8} & \textbf{55.5$\pm$2.0} \\
& E1-06 & FAU 4cl & Prosody & 65.0$\pm$1.2 & 54.0$\pm$2.2 \\
& E1-07 & IEMOCAP 4cl & Mean & 54.5$\pm$0.5 & 42.0$\pm$1.0 \\
& E1-08 & IEMOCAP 4cl & Self-Attn & \textbf{76.0$\pm$8.0} & \textbf{68.0$\pm$5.0} \\
& E1-09 & IEMOCAP 4cl & Prosody & 66.0$\pm$7.5 & 58.0$\pm$5.0 \\
\midrule
\multirow{6}{*}{E2}
& E2-01 & C-BESD Frozen & Self-Attn & \textbf{93.0$\pm$1.5} & 93.0$\pm$1.5 \\
& E2-01b& C-BESD Unfrozen & Self-Attn & 78.0$\pm$3.0 & 78.0$\pm$3.0 \\
& E2-02 & FAU Frozen & Self-Attn & \textbf{66.5$\pm$0.8} & 55.5$\pm$2.0 \\
& E2-02b& FAU Unfrozen & Self-Attn & 58.0$\pm$5.0 & 46.0$\pm$4.0 \\
& E2-03 & IEMOCAP Frozen & Self-Attn & 76.0$\pm$8.0 & 68.0$\pm$5.0 \\
& E2-03b& IEMOCAP Unfrozen & Self-Attn & \textbf{78.0$\pm$5.0} & 70.0$\pm$4.0 \\
\midrule
\multirow{18}{*}{E3$^\dagger$}
& E3-01 & C$\rightarrow$F & Mean & 18.0 & 22.0 \\
& E3-02 & C$\rightarrow$F & Self-Attn & 20.0 & 24.0 \\
& E3-03 & C$\rightarrow$F & Prosody & 19.5 & 23.5 \\
& E3-04 & C$\rightarrow$I & Mean & 28.0 & 26.0 \\
& E3-05 & C$\rightarrow$I & Self-Attn & 33.0 & 29.5 \\
& E3-06 & C$\rightarrow$I & Prosody & 30.0 & 28.0 \\
& $\cdots$ & $\cdots$ (12 remaining) & $\cdots$ & $\cdots$ & $\cdots$ \\
\midrule
\multirow{12}{*}{E4}
& E4-01 & C-BESD Clean & -- & 93.0$\pm$1.5 & 93.0$\pm$1.5 \\
& E4-02 & C-BESD +Adult & -- & 68.0$\pm$2.0 & 68.0$\pm$2.0 \\
& $\cdots$ & $\cdots$ (9 remaining) & $\cdots$ & $\cdots$ & $\cdots$ \\
\midrule
\multirow{9}{*}{E5}
& E5-01 & C-BESD LastL & Self-Attn & 80.0$\pm$2.0 & 80.0$\pm$2.0 \\
& E5-02 & C-BESD Best1 & Self-Attn & 84.0$\pm$1.5 & 84.0$\pm$1.5 \\
& E5-03 & C-BESD Weighted & Self-Attn & \textbf{93.0$\pm$1.5} & \textbf{93.0$\pm$1.5} \\
& $\cdots$ & $\cdots$ (6 remaining) & $\cdots$ & $\cdots$ & $\cdots$ \\
\midrule
\multirow{6}{*}{E6}
& E6-01 & C-BESD Mean+LastL & -- & 68.0$\pm$1.0 & 68.0$\pm$1.0 \\
& $\cdots$ & $\cdots$ (5 remaining) & $\cdots$ & $\cdots$ & $\cdots$ \\
\midrule
\multirow{6}{*}{E7}
& E7-01 & C$\rightarrow$F fine-tune & Self-Attn & 62.0$\pm$1.5 & 52.0$\pm$2.0 \\
& $\cdots$ & $\cdots$ (5 remaining) & $\cdots$ & $\cdots$ & $\cdots$ \\
\bottomrule
\end{tabular}
\end{table}
